\section{EVM rollup calldata-commitment costs: full derivation}
\label{app:eth-da}

This appendix gives the full, from-scratch derivation behind
Table~\ref{tab:eth-da-costs} and Section~\ref{sec:rollup-eth-results}:
every gas constant used, cited to its defining EIP, RSKIP, or ArbOS
release note; the per-option computation on each of the three chains;
and a fuller discussion of what the result does and does not establish.
The numbers below are computed in gas units (not ETH/RBTC/USD, which
float with market price and are not this paper's concern) against
current protocol constants as of this writing, with every non-protocol
approximation (implementation-dependent ``glue code'') flagged as such.

\subsection{A protocol fact the naive comparison misses}
\label{sec:rollup-eth-7623}

A gas-cost comparison based only on calldata's 16 gas/byte plus each
option's hashing cost misses a rule that binds precisely in this situation.
EIP-7623, live on Ethereum mainnet since the Pectra upgrade (May 2025), adds
a \emph{floor} to every transaction's gas cost \citep{eip7623}. Writing
$B$ for calldata bytes (worst case, all non-zero) and $E$ for all other
(``execution'') gas used -- hashing, storage writes, memory operations --
the total charged is
\begin{equation}
\mathrm{gasUsed} = 21000 + \max\big(16B + E,\ 40B\big),
\end{equation}
using EIP-7623's own parameters (\texttt{STANDARD\_TOKEN\_COST}$=4$,
\texttt{TOTAL\_COST\_FLOOR\_PER\_TOKEN}$=10$, and 4 ``tokens'' per non-zero
byte, giving $16B$ and $40B$ respectively). The floor $40B$ \emph{wins}
whenever $E < 24B$: whenever a transaction's real computation is cheap
relative to how much calldata it carries -- exactly the profile of a
DA-posting transaction, whose entire purpose is to carry bytes, not compute
on them. This means that for any of options B, C, D below, the specific
hashing algorithm's efficiency turns out not to matter once its execution
gas is (as we verify it is, in every case) comfortably under the
$24$-gas/calldata-byte threshold: the transaction is billed at the floor
regardless, and what remains decisive is only how many calldata bytes had
to be submitted in the first place. EIP-7623 is an Ethereum-mainnet
protocol rule, not a property every EVM-compatible chain shares by virtue
of being EVM-compatible -- Section~\ref{sec:rollup-rsk} below returns to
this once Rootstock is on the table, since Rootstock has no analogous
floor.

\subsection{The four options on Ethereum, computed}
\label{sec:rollup-eth-numbers}

Constants used, each unchanged since its introducing EIP and (for the
hash-related ones) not affected by any currently \emph{live} repricing --
two proposals that would raise them, EIP-2666 and EIP-7667, remain
Stagnant, and the draft Glamsterdam repricing set leaves compute opcodes and
precompiles unchanged (EIP-7904, listed in \citep{eip8007}); upcoming
changes that \emph{do} affect this comparison are discussed at the end of
this appendix:
\begin{itemize}
\item KECCAK256 opcode: $30 + 6\lceil \mathrm{len}/32\rceil$ gas.
\item SHA-256 precompile (\texttt{0x02}): $60 + 12\lceil
  \mathrm{len}/32\rceil$ gas.
\item BLAKE2F precompile (\texttt{0x09}): $1$ gas per round; the standard
  BLAKE2b compression uses 12 rounds, so 12 gas per call, each call
  processing one 128-byte message block against a 64-byte state
  \citep{eip152}.
\item Cold \texttt{SSTORE} of a zero storage slot to a non-zero value:
  \texttt{SSTORE\_SET\_GAS}$=20000$ \citep{eip2200} plus
  \texttt{COLD\_SLOAD\_COST}$=2100$ \citep{eip2929}, i.e.\ $22100$ gas per
  new 32-byte slot -- higher than a pre-2929 estimate of $20000$ alone.
\item A \texttt{STATICCALL} to a precompile is always ``warm'' post-Berlin,
  costing \texttt{WARM\_STORAGE\_READ\_COST}$=100$ gas
  \citep{eip2929} on top of the precompile's own listed price -- a cost the
  precompile gas figures above do \emph{not} include, and which matters
  once an option needs many separate precompile calls rather than one.
\end{itemize}

Let $P$ be the total original (unstructured) payload size in bytes, large
enough that fixed per-transaction terms (the $21000$ base, the small
constant terms in the hash formulas) amortize away.

\textbf{A -- persistent storage.} $B=P$; $E = (P/32)\times 22100 =
690.625P$. Since $E/B \gg 24$, the floor does not bind:
$\mathrm{gasUsed} \approx 16P + 690.625P = 706.625P$.

\textbf{B -- Keccak256 tree, discarded.} $B=P$. Using this paper's own
``same-convention'' tree (leaf $=H(\text{32-byte chunk})$, internal
$=H(\text{left}\Vert\text{right})$, matching Section~\ref{sec:benchmark}):
$N=P/32$ leaves, $2N-1\approx 2N$ calls, $E \approx 36N + 42N = 78N =
2.4375P$ (or $\approx 1.3125P$ if raw chunks are used directly as leaves,
skipping a separate leaf hash). Either way $E/B \ll 24$: the floor binds,
$\mathrm{gasUsed} = 40P$.

\textbf{C -- FlatMerkle, built off-chain, hashed via SHA-256.} The
structured file is $\approx 2P$ bytes ($2N{-}2\approx 2N$ items, $N=P/64$,
64 bytes each -- the same $\approx\!2\times$ blowup already measured in
Section~\ref{sec:benchmark}), all of it submitted as calldata: $B=2P$.
Hashing it in one SHA-256 precompile call costs $E \approx 0.375\times 2P =
0.75P$. Using BLAKE2F instead does not need a one-shot hash at all:
BLAKE2F only exposes the raw, one-block-at-a-time compression function
(unlike SHA-256's precompile, which hashes an arbitrary-length input in a
single call), so verifying the same structured file this way needs
$\approx 2N = P/64$ separate chained calls, each paying the $100$-gas
warm-call surcharge the one-shot SHA-256 call above never pays even once:
$E \approx (P/64)\times(100+12+\text{marshal}) \approx 2\text{--}3P$,
comparable to or worse than the SHA-256 figure, not better. Either way
$E \ll 24B = 48P$: the floor binds regardless of which precompile hashes
the file, at $\mathrm{gasUsed} = 40\times 2P = 80P$ -- exactly double
option B, purely because of the doubled calldata, not because of anything
about SHA-256 versus BLAKE2F.

\textbf{D -- FlatMerkle, built on-chain via BLAKE2F.} Only the original $P$
bytes are ever submitted as calldata; the structured file exists solely in
EVM memory. Treating BLAKE2b's block size the way
Section~\ref{sec:optimal-real} treats real SHA-256's (block size fixed,
every item -- data or checkpoint -- occupies one full block): $b=128$
bytes, $N=P/128$, $2N{-}1\approx 2N = P/64$ compression calls, each 12 gas
of pure precompile time ($E_{\text{pure}} = 0.1875P$). Accounting for the
warm-call surcharge and rough argument-marshaling cost each of the $P/64$
separate calls actually pays -- an engineering approximation, not a
protocol constant -- $E \approx (P/64)\times 135 \approx 2.109P$. Either
figure is far under the $24$-gas/byte threshold (an implementation would
need to be roughly \emph{ten times} less efficient than this rough estimate
before the floor stopped binding), so: $\mathrm{gasUsed} = 40P$.

\subsection{Rootstock: the same options, different protocol rules}
\label{sec:rollup-rsk}

Rootstock deliberately does not track Ethereum's gas schedule change for
change; RSKIP239 (Draft, unadopted, but authoritative on Rootstock's
\emph{current}, pre-proposal costs, which is what we need here) says so
directly: ``it is not our goal to match Ethereum's gas schedule. Indeed,
the current gas schedule in Ethereum has evolved even further away from
RSK's [schedule], post EIP-2929'' \citep{rskip239}. Concretely, as of this
writing Rootstock has adopted EIP-2028's calldata pricing (RSKIP400: 16
gas/non-zero byte, same as Ethereum \citep{rskip400}) and EIP-152's BLAKE2F
precompile (RSKIP153, same gas formula, 1 gas/round \citep{rskip153}), but
-- per RSKIP239's own ``current cost'' column, \texttt{SLOAD} still at 200
gas -- has \emph{not} adopted EIP-2929's Berlin repricing, and we found no
Rootstock proposal resembling EIP-7623's calldata floor at all. Two
consequences follow, in opposite directions:
\begin{itemize}
\item No calldata floor means Rootstock's transaction cost is the plain,
  unfloored sum $16B+E$ -- there is no $\max(\cdot,40B)$ waiting to
  absorb a cheap option's hashing-cost advantage the way it does on
  Ethereum.
\item No EIP-2929 means no warm-call discount for precompiles either:
  \texttt{STATICCALL} to a precompile still costs the flat, pre-Berlin
  \texttt{CALL} rate -- confirmed directly against the reference client's
  own gas-constants source, \texttt{CALL = 700} in \texttt{GasCost.java}
  \citep{rskjgascost}, with no cold/warm distinction defined anywhere in
  that file, cross-checking exactly against RSKIP239's account. This is
  the same flat 700 that EIP-2929 itself replaced with the 100/2600 split
  on Ethereum \citep{eip2929}. This makes \emph{repeated} precompile
  calls, not just hashing throughput, expensive on Rootstock in a way
  they are not on Ethereum.
\end{itemize}
\texttt{SSTORE} of a fresh 32-byte slot, with no cold-slot surcharge to
add, stays at the plain \texttt{SET\_SSTORE = 20000}
\citep{eip2200,rskjgascost} -- $625$ gas/byte, exactly the figure a
pre-EIP-2929 estimate would give; the same source also confirms
\texttt{SHA3 = 30}, \texttt{SHA3\_WORD = 6} and \texttt{SHA256 = 60},
\texttt{SHA256\_WORD = 12}, identical to the Ethereum values
Section~\ref{sec:rollup-eth-numbers} uses.

Recomputing all four options against these rules ($B$, $E$, $P$ as
before, no floor -- $\mathrm{gasUsed}=16B+E$ directly):

\textbf{A.} $E=(P/32)\times 20000=625P$; total $=16P+625P=641P$.

\textbf{B.} Unchanged hashing cost ($1.3$--$2.4P$, Keccak256 is a native
opcode on both chains, needing no call at all); total $=16P +
(1.3\text{--}2.4)P \approx 17.3\text{--}18.4P$.

\textbf{C.} $B=2P$ as before. One-shot SHA-256: $E\approx 0.75P$
(the flat per-call overhead is paid once, negligible for large $P$);
total $\approx 32.75P$. Chained BLAKE2F: $E\approx (P/64)\times
(700+12+\text{marshal})\approx 11.5P$; total $\approx 43.5P$ -- here,
without a floor to absorb it, using BLAKE2F instead of SHA-256 makes this
option \emph{more} expensive, not merely no cheaper, since Rootstock's
flat 700-gas call cost is paid $P/64$ times instead of once.

\textbf{D.} $B=P$, $E\approx (P/64)\times(700+12+\text{marshal})\approx
11.5P$; total $\approx 27.5P$.

Ranked cheapest to most expensive on Rootstock: B ($\approx\!17.3$--$18.4$)
$<$ D ($\approx\!27.5$) $<$ C ($\approx\!32.75$--$43.5$) $\ll$ A ($641$).
This is not the ranking we expected going in, and it is worth stating
plainly rather than smoothing over: on Rootstock, option D does
\emph{not} beat option B -- removing the calldata floor does let raw
hashing efficiency matter again, but Rootstock's absence of a warm-call
discount taxes option D's $P/64$ separate BLAKE2F calls more than it
costs option B's zero calls (a native opcode, not a precompile) or even
option C's single call. What Rootstock's numbers \emph{do} show is that
D is a clear, large improvement over C (off-chain construction) and over
A (persistent storage) on any EVM chain we have checked, and that which
of B or D is cheapest is not a fact about FlatMerkle or about the EVM in
general -- it is a fact about whether the specific chain discounts
repeated precompile calls the way Ethereum has since Berlin, which a
deployer needs to check for their actual target chain rather than assume.

\subsection{Arbitrum: the missing combination}
\label{sec:rollup-arb}

Ethereum has both a calldata floor and a warm-call discount (the discount
just never gets to matter, since the floor already dominates); Rootstock
has neither. The remaining combination -- a warm-call discount
\emph{without} a calldata floor -- is not hypothetical: it is Arbitrum's
current, documented configuration. Arbitrum One and Nova run Nitro, a
modified Geth core that tracks Ethereum's Istanbul/Berlin/London opcode
and precompile pricing, including EIP-2929's cold/warm split (2100/100 for
\texttt{SLOAD}, the same split applied to precompile
\texttt{STATICCALL}s); its own documentation, however, states plainly that
EIP-7623 was \emph{not} enabled as part of ArbOS 40 Callisto (Arbitrum's
Pectra-alignment upgrade, activated 2025) and remains unenabled as of
ArbOS 51 Dia (activated January 8, 2026, current as of this writing),
``because block size variance is less of a concern on Arbitrum chains''
\citep{arbos40,arbos51}. This is precisely the case relevant to an
Arbitrum Orbit L3 that wants to use Arbitrum One itself (rather than
Ethereum L1) as its settlement/DA layer -- an L2-native analogue of the
Bitcoin rollup setting of Section~\ref{sec:rollup-da}, with Arbitrum
playing the role of the underlying, published-to chain.

Recomputing with Ethereum's own hashing/call constants (Arbitrum inherits
them unchanged) but \emph{no} EIP-7623 floor -- $\mathrm{gasUsed}=16B+E$,
as on Rootstock, but with $100$-gas warm precompile calls, as on Ethereum:

\textbf{A.} Same as Ethereum: $706.625P$.

\textbf{B.} Unchanged: $16P + (1.3\text{--}2.4)P \approx
17.3\text{--}18.4P$.

\textbf{C.} $B=2P$; one-shot SHA-256 $\approx 32.75P$; chained BLAKE2F
$E\approx(P/64)\times(100+12+\text{marshal})\approx 2.1P$, total
$\approx 34.1P$ -- still no cheaper than one-shot SHA-256, for the same
reason as before (many calls, even discounted ones, against one).

\textbf{D.} $B=P$, $E\approx(P/64)\times 135\approx 2.1P$ (realistic) or
$0.1875P$ (pure precompile); total $\approx 18.1P$, or as low as
$16.1875P$ under the optimistic estimate.

Here, unlike Rootstock, option D is at least competitive with option B
and plausibly a little cheaper ($16.2$--$18.1P$ against
$17.3$--$18.4P$), rather than clearly worse -- though the two ranges
overlap enough, given how approximate the ``glue code'' estimate is, that
we would not call this a decisive win either. What removing the floor
while \emph{keeping} the discount does establish is directional: it moves
D from a clear loss (Rootstock) to a close contest (Arbitrum), consistent
with the mechanism identified above -- the discount is what lets many
cheap calls compete with zero calls at all, and only shows up once
nothing else swamps it.

\subsection{Two narrower points, and the general lesson}
\label{sec:rollup-eth-conclusion}

FlatMerkle's case here rests on two narrower points, not on beating a
plain Merkle tree's cost on every chain. First, the SHA-256 precompile's
one-shot, arbitrary-length interface cannot build a structured file
incrementally with a custom starting state the way BLAKE2F's raw,
resumable compression-function interface can (this is exactly option C's
cost penalty above); BLAKE2F is, as far as we are aware, the only
precompile available on any of the three chains that exposes what
FlatMerkle's construction actually needs, so a rollup that wants the
structured file built entirely on-chain -- with no dependence on an
off-chain party supplying it correctly, and no fraud-proof step to check
that they did -- has option D as close to its only route to that
property, not merely its cheapest one. Second, a rollup that also
publishes a redundant or backup commitment via the Bitcoin annex
application of Sections~\ref{sec:bitcoin-annex}--\ref{sec:rollup-da} can
reuse a single flat, linear proof format across all four venues (Bitcoin,
Ethereum, Rootstock, Arbitrum); a Keccak Merkle tree is not that format on
any of them, and would need its own, separate verifier per venue.

More generally: whether option B or D is cheapest is not a fact about
FlatMerkle, nor about ``the EVM'' as a single target -- it is a fact
about two \emph{separable} protocol design choices a given chain may or
may not have made: a calldata floor (EIP-7623-style, which erases the gap
between B and D by making both floor-bound) and a discount for repeated
precompile calls (EIP-2929-style, which is what lets D's many small calls
compete with B's zero calls at all once there is no floor to erase the
difference). Ethereum currently has the first, and the second without it
mattering (it applies, but the floor already dominates); Rootstock has
neither; Arbitrum has the second without the first, which is the one
combination among the three where option D is not clearly worse than
option B. A rollup deployer targeting a different EVM-compatible chain
should check both properties directly rather than assume any of these
three answers transfers.

Two further things this analysis does not cover, both flagged rather than
resolved. First, the exact per-call ``glue code'' cost: D's $\approx 135$
gas/call on Ethereum and Arbitrum and $\approx 737$ on Rootstock are
estimates for a reasonably careful Solidity or Yul implementation, not
measured or protocol-fixed numbers (on Ethereum the margin to the floor
threshold is wide enough that this does not change the conclusion).
Second, every constant used here is a protocol \emph{parameter}, and
several are about to move. Ethereum's next fork, Glamsterdam, was
activated on the Sepolia testnet in October 2026 and has no mainnet date
as of this writing; its draft repricing set \citep{eip8007} raises the
calldata floor from $10/40$ to $64/64$ gas per byte (EIP-7976), which
would make the floor bind whenever $E < 48B$, putting options B and D at
$64P$ and option C at $128P$ -- the same ranking at higher absolute cost --
and raises the cost of creating a storage slot to $97{,}920$ state gas
(EIP-8037), making option A several times more expensive still. The
draft's numbers are marked preliminary. Further out, EIP-7667, if ever
adopted, would raise \texttt{BLAKE2\_GFROUND} from 1 to 10 and the
SHA-256/Keccak costs by roughly 5--10$\times$, moving options B, C and D
toward the floor threshold; RSKIP239, if adopted, would not touch the
costs of Section~\ref{sec:rollup-rsk}, but shows that Rootstock's schedule
is also a moving target. The tables are the correct computation against
the rules in force at the time of writing, not a permanent ranking.

As illustrative, non-binding capacity figures: since Fusaka (December
2025), EIP-7825 caps every Ethereum transaction at $2^{24} =
16{,}777{,}216$ gas \citep{eip7825}, so a single floor-bound transaction
(options B or D) carries at most $(2^{24} - 21{,}000)/40 \approx 419{,}000$
bytes, and filling a $\approx\!60$-million-gas block, about $1.5$ million
bytes, takes several transactions. At Rootstock's $6.8$ million gas block
limit \citep{rskgas}, option B's $\approx\!17.3$--$18.4$ gas/byte allows
roughly $370$--$390$ thousand bytes. At Arbitrum's $32$ million gas
per-transaction cap (ArbOS 51 Dia, child-chain execution gas only,
excluding any separate parent-chain data fee) \citep{arbos51}, option D's
$\approx\!16.2$--$18.1$ gas/byte allows roughly $1.8$--$2.0$ million bytes
in a single transaction. All three figures ignore every other use of the
same block's or transaction's gas.
